% ===================== 摘要 =====================
% 规则（详见技能 references/paper-writing.md）：
%   总字数 <= 900 字且必须一页；开头段 <= 30 字；结尾段 <= 20 字。
%   每问一段，按问题类型弹性分配字数：建模求解型 30-35%、数据分析型 20-25%、建议总结型 10-15%。
%   每问必须出现具体数值结果；关键词 5 个左右。
\begin{abstract}

% 开头段（<=30 字）：
...是...领域的重要研究课题。本文基于...数据，围绕...展开系统研究。

% 逐问摘要（每问选一种类型模板替换【】）：
% ① 有数据型：
\textbf{针对问题一，}本问是【类型】问题，需【要求】。首先对附件数据预处理，提取【特征】，【清洗/标准化】。接着进行【变量筛选/降维/参数选择】。然后建立【模型】进行【求解过程】。最后通过【对比/验证】，结果表明\textbf{【关键数值结果】}，模型有效地【评价】。

% ② 无数据（机理）型：
\textbf{针对问题二，}本问是【类型】问题，需【要求】。首先分析【机理/物理过程/目标约束】，建立【方程/优化模型】。接着设定【参数/初始条件】。然后通过【数值方法/算法】求解，得到【结果】。结果表明\textbf{【关键数值结果】}，模型有效地【评价】。

% ③ 建议总结型（一笔带过）：
\textbf{针对问题三，}本问是建议总结问题。基于前文分析，从【角度1】和【角度2】出发给出【N】条建议，核心结论为\textbf{【要点】}。

% 结尾段（<=20 字）：
本文模型性能优良、结论可靠，可为【领域】提供决策参考。

\keywords{关键词1；关键词2；关键词3；关键词4；关键词5}
\label{abstract:end}
\end{abstract}
